oz-policy-builder is ahead of Perch in three respects, which a reader choosing a
design today should weigh first.

\paragraph{Maturity.} oz-policy-builder pins one interpreter instance on testnet
(\srcpath{packages/policy-synth/src/config.ts:57-80}). Release \fact{opb-release} is
the code submitted for an external audit, which started on
\fact{opb-audit-start} and is \fact{opb-audit-status}, with the OpenZeppelin
account and policies outside its scope (\srcpath{README.md:57-59};
\srcpath{docs/audit-overview.md:76-85}). Its off-chain packages were hardened after
an external penetration test, in pull request \fact{opb-pentest-pr} (branch
\code{\fact{opb-pentest-branch}}), merged at the pin. The interpreter has
\fact{opb-interpreter-tests} tests, \fact{opb-tests-against-oz08} of them against
OpenZeppelin's \fact{oz-v080-tag} account (\srcpath{docs/audit-overview.md:15-20}); a
corpus test takes \fact{opb-corpus-calls} real mainnet transaction envelopes, of
\fact{opb-corpus-shapes} shapes, through decoding, synthesis and the interpreter
(\srcpath{packages/policy-synth/src/corpus.test.ts:1-3}); mutation testing of one
path (mutants: \fact{opb-mutants}); and a STRIDE register gives each threat a
control and a status. Perch's grant design is an open pull-request stack ending at
\fact{perch-pr-stack-head}, runs on testnet only as a preview, and is unaudited.

\paragraph{Authoring.} oz-policy-builder starts from what an app did. Its
synthesiser records a transaction, or takes a declared Intent, and writes the
least predicate that permits it; verifies it by dry-running the deployed
interpreter's own \code{check}, so one evaluator both verifies and enforces;
drops each conjunct in turn to show that each is needed (minimality); renders a
review; and returns unsigned install bytes
(\srcpath{packages/policy-synth/README.md}, lines 58 and 142). Its overlap scan reads every
live rule of the account and flags, among other cases, a rule a shared signer
could bypass through a looser neighbour (\srcpath{policy-synth/README.md:171-176}).
The same \fact{opb-tools-count} tools run from a command line and from an MCP
server, so an agent can drive them, and the packages are published at
\fact{opb-packages-version}. Perch starts from what the owner grants: a wallet
declares a grant document, which the account's compiler validates and the
account lowers against its own authority. Perch has a document builder and
offline checks, but no recorder, synthesiser or review renderer.

\paragraph{Expressiveness and ecosystem fit.} The interpreter's language has
\code{and} and \code{or}, six comparisons, \code{in} and \code{nin}, \code{all},
\code{any} and \code{sum} over vectors, scaled arithmetic, the ledger time
\code{now}, nested \code{call\_path} and \code{read}
(\srcpath{policy-interpreter/src/dsl.rs:84-158}). An Intent can state a clock window
(\code{until}, \code{from}), a least swap output relative to another argument
(\code{minOut}), and conditions on every, some or the total of a vector's elements
(\code{each}, \code{some}, \code{total}) (\srcpath{intent.ts:18-49}). A rule can read
up to \fact{opb-interpreter-max-reads} values from other contracts. Multi-signer
rules use OpenZeppelin's threshold policies. A refusal reports a specific code
for each reason (\srcpath{policy-interpreter/src/storage.rs:27-86}), where Perch
reports \code{StaleScope} (\fact{auth-error-stale-scope}) or
\code{RuleEvalRefused} (\fact{auth-error-rule-eval-refused}). And OpenZeppelin
\fact{oz-v080-tag} binds the selected rule IDs for both signer kinds
(\srcpath{storage.rs}, lines 157--175 and 396--412, in \code{stellar-contracts}), where Perch's
delegated signers do not bind them.
